\documentclass[11pt]{article}
\usepackage{amsmath,amssymb,amsthm,hyperref}
\newtheorem{theorem}{Theorem}
\newtheorem{lemma}{Lemma}
\title{Degree-aware dyadic readout: exact certificates and a precision obstruction}
\author{Research checkpoint 008}
\date{31 October 2026}
\begin{document}
\maketitle

\paragraph{Scope.}
We optimize an exact linear identity on \emph{all rational polynomials}
of a given degree bound. This is not the problem of recovering only a
constant coefficient modulo two on an integer-coefficient or
tensor-realizable subclass. It is not an ECDLP complexity lower bound.
We find real finite improvements over the standard full Hermite
functional, certify their optimality, and prove that optimizing this
functional alone cannot improve the asymptotic geometric precision
scale on the consecutive grid. We also give corresponding lower bounds
for the two balanced-cluster node sets studied in checkpoint 007.
No complete contraction recurrence or complete ECDLP exponent improves.

\section{A finite, constructive optimum}
Let distinct integer centers $\tau$ have jet lengths $q_\tau$ and local
scales $\ell_\tau\ge0$. The rows indexed by $(\tau,j)$ of the integer
matrix $A$ are
\[
 A_{(\tau,j),k}=\binom{k}{j}\tau^{k-j}2^{\ell_\tau j},
 \qquad 0\le k\le d,\quad 0\le j<q_\tau,
\]
with zero entries when $j>k$. Assume $\sum q_\tau\ge d+1$.
The common objective is the least integer $L\ge0$ such that
$wA=e_0$ with rational weights satisfying $v_2(w_r)\ge-L$.

Choose initial jet prefixes of sizes $m_\tau\le q_\tau$ totaling $d+1$.
Their square confluent Vandermonde determinant, up to sign, is
\[
 2^{\sum_\tau\ell_\tau\binom{m_\tau}{2}}
 \prod_{\tau<s}(s-\tau)^{m_\tau m_s}.
\]
There are no derivative factorials: these are Taylor coefficients.
For example, coalesce $m_\tau$ ordinary Vandermonde points at each
center, divide out the within-center Vandermonde factors, and compare
the first nonzero coefficient; the remaining rows are precisely the
Hasse derivatives. Let $C$ be this determinant's valuation.
Multiplying row $r$ by $2^{p_r}$ adds $\sum p_r$ over the minor to $C$.

The determinant bound $C$ can be very conservative for the working
precision. Here is a smaller proved bound for inverse-minor entries.
Let $G$ be the maximum pair-difference valuation among used centers
(zero for a single center), and
$S_\tau=\sum_{s\ne\tau}m_s v_2(\tau-s)$.
For the minor column dual to jet $(\tau,j)$, an inverse-entry loss bound is
\[
 B_{\tau,j}=S_\tau+(m_\tau-1-j)G+\ell_\tau j+p_{\tau,j}.
\]
Indeed its unscaled cardinal polynomial is
$(T-\tau)^j E_\tau(T)$ times the inverse Taylor series of $E_\tau$
at $\tau$, truncated after degree $m_\tau-1-j$, where
$E_\tau(T)=\prod_{s\ne\tau}((T-s)/(\tau-s))^{m_s}$.
The polynomial $E_\tau$ loses at most $S_\tau$ bits; its reciprocal
Taylor coefficient of order $k$ loses at most $kG$ bits. Divide by
the local row scale $2^{\ell_\tau j+p_{\tau,j}}$.
This proves the bound for every coefficient, including unequal scales
and additional row powers. Put $B=\min(C,\max B_{\tau,j})$.

Work modulo $2^K$, with $K=B+2$. Invertible row and column operations
using minimum-valuation pivots give $UAV=D$ modulo $2^K$, where the
nonzero diagonal entries are $2^{\nu_i}$. The known minor guarantees
$\sum\nu_i\le C$ and $\max\nu_i\le B$, so no invariant factor is
hidden by the cutoff.
Put $t=e_0V$. Then
\begin{equation}
 L=\max\bigl(0,\max_i\{\nu_i-v_2(t_i)\}\bigr),
 \label{eq:optimum}
\end{equation}
where a zero residue has valuation at least $K$.

\begin{proof}[Certification of \eqref{eq:optimum}]
The inverse-minor bound implies that the row module over $\mathbb Z_2$
contains $2^B\mathbb Z_2^{d+1}$. Thus membership modulo $2^K$ is enough
for exact membership. The diagonal matrix gives the stated condition.

For an exact rational upper certificate, solve
$xD=2^Lt$ modulo $2^K$, and let $\lambda=xU$ be an integer lift.
The residual $2^Le_0-\lambda A$ is divisible by $2^K$.
Correct it using the chosen square minor. Its inverse has valuation
at least $-B$ entrywise, so the correction has valuation at least
$K-B>0$. The corrected coefficients divided by $2^L$ are exact
rational weights with the required denominator bound.

If $L>0$, choose a column attaining the maximum in
\eqref{eq:optimum}, and lift that column of $V$ to an integer
coefficient vector $f$. Its constant coefficient has valuation $t_0$,
and its minimum jet valuation is exactly $t_0+L$. Weights of loss at
most $L-1$ would reconstruct a number of valuation at least $t_0+1$,
contradicting $f(0)$. This is an integer-polynomial lower certificate
that can be checked without repeating the diagonalization.
\end{proof}

As another check, augment the row module by $e_0$. The decrease in the
sum of invariant-factor valuations is $L$: the new quotient is cyclic,
generated by $e_0$, with annihilator $2^L$.

For unequal nonnegative budgets $b_\tau$, let $B=\max b_\tau$ and
multiply row $(\tau,j)$ by $2^{B-b_\tau}$. The original budgets are
feasible exactly when the common optimum of this scaled matrix is
at most $B$. Binary search therefore optimizes the geometric profile
\begin{equation}
 b_\tau=\left\lfloor\frac{\lambda}{2^{v_2(\tau)}}\right\rfloor
 \quad(\tau\ne0).
 \label{eq:profile}
\end{equation}
The saved result includes exact upper weights and a polynomial proving
infeasibility at $\lambda-1$.

The finite solver uses polynomially many bit operations in the explicit
matrix dimensions and bit heights, and polynomial storage: the modular
matrix and transforms have $O(mn+m^2+n^2)$ entries of $K$ bits.
The exact correction has polynomial bit length by determinant bounds.
These costs must still be paid in any proposed contractor. The finite
solver is not a constant-cost oracle for a succinct exponential matrix.

\section{A lower certificate valid for any available jet lengths}
Let $\mathcal N$ be a finite set of nonzero integer centers. Suppose
$\ell\ge\max_{s\ne\tau}v_2(\tau-s)$ and $\ell\ge v_2(\tau)$.
Choose positive integer root orders $a_\tau$ with
$\sum a_\tau\le d$, and set
\[
 f(T)=\prod_{\tau\in\mathcal N}(T-\tau)^{a_\tau}.
\]
At center $\tau$ the first nonzero scaled jet has order $a_\tau$.
All scaled jets have valuation at least
\[
 \ell a_\tau+\sum_{s\ne\tau}a_s v_2(\tau-s),
\]
and equality holds at that first nonzero jet. This follows by expanding
each factor: its scaled linear coefficient has at least the valuation
of its nonzero constant coefficient. Define
\begin{equation}
 D_\tau=(\ell-v_2(\tau))a_\tau+
 \sum_{s\ne\tau}a_s\bigl(v_2(\tau-s)-v_2(s)\bigr).
 \label{eq:dual}
\end{equation}
Every available jet then has valuation at least $v_2(f(0))+D_\tau$.
If its order is less than $a_\tau$, it is zero instead.
Consequently every exact linear readout obeying \eqref{eq:profile}
must satisfy
\begin{equation}
 \lambda\ge\min_{\tau\in\mathcal N}2^{v_2(\tau)}D_\tau.
 \label{eq:lower}
\end{equation}
This bound remains valid if nodes are omitted or more jets are made
available. It applies whenever an exact readout exists; a collection
whose jets all vanish on $f$ cannot reconstruct this polynomial at all.

\section{Consecutive coordinates}
Let $L=2^a$, $1\le u\le a$, and retain centers
$-L\le\tau<L$ with $v_2(\tau)<u$. Take $\ell=a+1$ and degree
$d=M(2L-1)$. Put
\[
 J=2L(u-1+2^{-u}),\qquad h=\lfloor d/J\rfloor,\qquad
 a_\tau=h(u-v_2(\tau)).
\]
The witness degree is $hJ\le d$. For a center at level $i$, its
cross-center contribution in \eqref{eq:dual} is
\[
 H_i=h\frac{L}{2^i}
       \sum_{j=1}^{a-i}\frac{j}{2^j}\min(j,u-i)
 \ge\frac{hL}{2^{i+1}}.
\]
The equality follows by counting centers in each dyadic distance
class; it is also the negative valuation of the full Hermite factor
$\prod_{s\ne\tau}(-s/(\tau-s))^{a_s}$.
Since the local term in \eqref{eq:dual} is nonnegative,
\[
 \lambda\ge hL/2.
\]
If $M\ge4u$, then $d\ge ML$, $J\le2Lu$, and
$h\ge M/(4u)$. Therefore
\begin{theorem}
For these consecutive centers and local scales, every exact linear
readout on all polynomials of degree at most $M(2L-1)$ with the
geometric precision profile satisfies
\[
 \lambda\ge\frac{ML}{8u}\qquad(M\ge4u).
\]
This holds regardless of the available jet multiplicities.
\end{theorem}

For comparison, when $u\ge2$, the standard choice
$q_i=\lceil M/(u-1)\rceil(u-i)$ gives a matching $O(ML/u)$ upper bound.
Write $h'=\lceil M/(u-1)\rceil$. The cross-center loss satisfies
$2^iH_i\le6h'L$ using $\sum_{j\ge1}j^2/2^j=6$.
The local contribution is at most $\tfrac32h'L$, because
$2^{-j}j(j+1)\le3/2$. These bounds apply to every Hermite coefficient:
the local inverse series is integral, and the highest-jet coefficient
attains the largest possible loss. Thus $\lambda\le8h'L\le24ML/u$
under $M\ge4u$. Degree-aware linear readout changes constants, but
cannot change the order $ML/u$ of this precision profile. Root output
precision, normalization and recursive losses remain separate costs.

\section{Balanced geometric clusters}
Use the notation of checkpoint 007:
\[
 B=2^u,\quad c=(n-3)(n-4)+1,\quad e=2(n^2-5n+9),\quad n\ge4.
\]
The difference grid has $Be-c$ points, the deformation degree is
$d=M(Be-c-1)$, and $\ell=u+n-1$ exceeds every pair-difference
valuation. Retain levels $i<v\le u$.

\subsection*{All coarse nodes}
Set $J=eB(v-1+2^{-v})$, $h=\lfloor d/J\rfloor$, and
$a_\tau=h(v-i)$. At a node $\tau=Bt+r$ with $i=v_2(\tau)$, put
$k=u-i$ and $z=v-i$. The exact dyadic distance count from checkpoint
007 gives
\[
 H_\tau=h\left[e\mathcal A(k,z)+
 z((e-1)k+V_{E_n}(t))\right],
\quad
 \mathcal A(k,z)=3\cdot2^k-2^{k-z}(z+3)-z(k+1).
\]
Here $V_{E_n}(t)\ge0$ and $e\ge2$.
The function $\mathcal A(k,z)$ is nondecreasing in $z\le k$, and
$\mathcal A(k,1)=2^k-k-1$. Hence
\[
 \frac{H_\tau}{he}\ge2^k-k-1+\tfrac12zk
 \ge2^k-k/2-1\ge2^k/4.
\]
The last inequality holds for $k\ge1$ (equality at $k=1$).
Thus \eqref{eq:lower} gives $\lambda\ge heB/4$.
Since $c+1\le e$, $B\ge2$, $d\ge MeB/2$, and $J\le eBv$,
$h\ge M/(4v)$ for $M\ge4v$. We obtain
\[
 \boxed{\lambda\ge\frac{MeB}{16v}
                   \ge\frac{Mn^2B}{32v}}.
\]

\subsection*{Positive central nodes only}
Now $1\le\tau<B$, $i<v$, and $J=B(v-1+2^{-v})$.
Use $h=\lfloor d/J\rfloor$ and the same orders.
The cross-center contribution is $h\mathcal A(k,z)$.
Since $\ell-i\ge k$, the entire loss obeys
\[
 D_\tau\ge h(\mathcal A(k,z)+kz)
             \ge h(2^k-1)\ge h2^{k-1}.
\]
Therefore $\lambda\ge hB/2$. For $Me\ge4v$,
$h\ge Me/(4v)$, giving
\[
 \boxed{\lambda\ge\frac{MeB}{8v}
                   \ge\frac{Mn^2B}{16v}}.
\]

These cluster bounds close the specific loophole of achieving a
smaller \emph{geometric operator-precision scale} by changing the
linear functional. They do not prove that every node must be computed,
that every feasible mode contributes, or that a different precision
profile has the same contraction cost. A new complete candidate-count
or runtime theorem would need those additional arguments. Integrality
of the underlying packet tensors is also a separate requirement;
the algebraic readout bound does not waive it.

\section{Reproducibility and remaining direction}
The launcher retains exact rational upper identities, independent
integer-polynomial lower certificates, unique full-Hermite controls,
confluent determinant checks, augmented-module index checks, and
negative controls. An independent checker rebuilds jet coefficients
by polynomial multiplication, without using the optimizer's matrix
builder or diagonalization. Small instances expand the obstruction
polynomial and check every coefficient of every local Taylor series;
large instances only check closed formulas and are labeled accordingly.

Finite degree fixtures are not executions of a complete C38 network.
Cases whose mass is smaller than one packet are labeled polynomial
fixtures. Any future gain must come from something beyond this
degree-only geometric precision change: e.g. a justified constraint on
the actual deformation space, a different cost profile with a complete
recurrence, or a representation that avoids unrestricted mode lists.
The general square-root ECDLP target remains active and unachieved.
\end{document}
